% ============================================================
%  MÓDULO — Instalação, Uso e Operação do Core
% ============================================================
\chapter{Instalação, Uso e Operação}

\roteirodidatico
{Instalar significa preparar o ambiente para executar o programa. Um comando que termina sem erro é um primeiro sinal, mas ainda não demonstra que serviços, credenciais e modelos necessários estão disponíveis.}
{\emph{Dependência} é um programa ou pacote exigido; \emph{ambiente virtual} isola pacotes Python; \emph{provedor} entrega inferência de modelo; \emph{diagnóstico} verifica condições locais.}
{Siga a instalação por etapas, confira a saída de cada comando e diferencie componentes obrigatórios dos opcionais. Use o diagnóstico para identificar qual dependência está faltando.}
{Reprodutibilidade requer registrar sistema, versões, perfil instalado, configuração e provedores habilitados; “instalado” não significa “conectado” nem “validado”.}
{Que pré-requisito foi satisfeito, por qual evidência local, e qual dependência externa ainda pode impedir a tarefa?}

\casoguiado{o diagnóstico responde parcialmente}
{Se o diagnóstico informa que as especificações carregaram, mas uma credencial de provedor não foi encontrada, as duas observações devem permanecer distintas: o mecanismo de leitura está disponível; a inferência que depende daquela credencial talvez não esteja. O passo seguinte é verificar a configuração e repetir apenas a checagem pertinente. Não se conclui que todo o ambiente está quebrado nem que o sistema está plenamente operacional.}

\section{O que você precisa}

\elemento{Ambiente mínimo}{\metainfo{ID}{U-01} \hfill
\metainfo{Rota}{\pth{requirements*.txt}}}

\begin{descricao}
Para operar o Core você precisa de: Python 3 (ambiente virtual recomendado),
ferramentas de linha de comando do sistema (git, conforme o caso), e, se quiser
inferência local, o servidor \textbf{Ollama} com modelos (veja operação local
adiante). Modelos da nuvem (OpenCode Zen, OpenAI, etc.) exigem credenciais
próprias.
\end{descricao}

\begin{metodo}
\textbf{Passo a passo de instalação:}
\begin{enumerate}[leftmargin=1.4em,itemsep=1pt]
  \item Clone/atualize o repositório e crie o venv:
  \begin{itemize}[leftmargin=1.3em,itemsep=0pt]
    \item \texttt{make\ install}
  \end{itemize}
  \item Instale as dependências base e (opcional) as de desenvolvimento/ciência:
  \begin{itemize}[leftmargin=1.3em,itemsep=0pt]
    \item \texttt{make\ setup}
    \item \texttt{make\ setup-dev} (desenvolvedores)
    \item \texttt{make\ setup-science} (científico)
  \end{itemize}
  \item Valide a saúde do ambiente (regra do \pth{AGENTS.md}):
  \begin{itemize}[leftmargin=1.3em,itemsep=0pt]
    \item \texttt{python3\ -m\ marceloclaro.cli\ doctor}
  \end{itemize}
  \item Se houver um instalador específico da plataforma, siga as instruções
        correspondentes em \pth{common/} e consulte \pth{specs/} quando precisar
        da especificação.
\end{enumerate}
\end{metodo}

\begin{funcao}
Garantir que o Core consiga: carregar specs, ler o registro de evolução,
acessar a memória metacognitiva e despachar agentes — exatamente o que o
\texttt{doctor} confere.
\end{funcao}

\begin{niveis}
\textbf{Intuição:} como instalar um ``programa com muitos ajudantes'' — uma pasta com
regras e um comando de cheque médico (\texttt{doctor}).\\
\textbf{Implementação:} venv isolado + requirements por perfil (base/dev/scientific),
instaladores por SO e unit para systemd.\\
\textbf{Formalização:} o \texttt{doctor} é um oráculo de health-checks (20 checks na última
configuração registrada) ligando specs $\leftrightarrow$ código $\leftrightarrow$
memória; é a primeira etapa do protocolo de auditoria.
\end{niveis}

\clearpage
\section{Como usar no dia a dia}

\elemento{Comandos principais}{\metainfo{ID}{U-02} \hfill
\metainfo{Rota}{\pth{MANUAL.md}}}

\begin{funcao}
Operar o ecossistema sem editar arquivos de configuração manualmente.
\end{funcao}

\begin{processo}
\textbf{Menu de operação} (visão consolidada do \pth{CLI}):
\begin{enumerate}[leftmargin=1.4em,itemsep=1pt]
  \item \texttt{python3\ -m\ marceloclaro.cli\ menu} — menu principal.
  \item \texttt{python3\ -m\ marceloclaro.cli\ doctor} — diagnóstico de saúde (rode
        \emph{antes} de qualquer edição).
  \item \texttt{python3\ -m\ marceloclaro.cli\ helpdesk} — ajuda guiada em linguagem
        simples.
  \item \pth{opencode.json},
        catálogo de agentes e comandos customizados (após editar cards em
         \pth{skills/}).
  \item Skills instaladas em \pth{.opencode/skills/} — pesquisa científica,
        médico virtual supremo, nano-orquestração, deploy GitHub Pages, MiroFish,
        prestação de contas, entre outras.
\end{enumerate}
\end{processo}

\begin{metodo}
\textbf{Fluxo recomendado para uma tarefa:}
\begin{enumerate}[leftmargin=1.4em,itemsep=1pt]
  \item \texttt{python3\ -m\ marceloclaro.cli\ doctor} para conferir saúde.
  \item Invocar o orquestrador primário \pth{marceloclaro} (ponto de entrada único
        para tarefas multipasso).
  \item O orquestrador percebe, delega pelo blackboard e exige spec+testes.
  \item Ao terminar, confira o registro de evolução (\pth{evolution/}).
\end{enumerate}
\end{metodo}

\begin{limite}
Permissões padrão do Core: \texttt{edit: ask} e \texttt{bash: ask} — o operador
aprova cada ação; não há execução silenciosa.
\end{limite}

\section{Operação com modelos locais (Ollama)}

\elemento{Perfil otimizado de inferência local (SPEC-935-R607)}{\metainfo{ID}{U-03} \hfill
\metainfo{Rota}{\pth{specs/SPEC-935-R607.md}}}

\begin{descricao}
O Core pode rodar com modelos locais via Ollama. A otimização R607 mediu o custo de
contexto e definiu um \emph{perfil mínimo} de configuração, mais veloz, para uso
local.
\end{descricao}

\begin{resultado}
\textbf{Medições verificadas (R607, no host de teste):}\par
\smallskip
{\footnotesize
\begin{tabularx}{\textwidth}{Xcc}
\toprule
Configuração & Tamanho do prompt & Tempo por turno \\
\midrule
Perfil completo (catálogo inteiro) & $\approx$ 21.136 tokens & 6 a 18 min (timeouts de 500) \\
Perfil mínimo local & 2.051 a 6.712 tokens & 1 min 32 s a 5 min 01 s \\
\bottomrule
\end{tabularx}\par}
\textbf{Ajustes-chave:} \pth{num_ctx=24576} nos modelos
afinados \pth{gemma4:e4b-tuned}; no systemd,
\pth{OLLAMA_MAX_LOADED_MODELS=1}. O valor
``16,20~tok/s'' de uma medição anterior foi \textbf{revogado} (artefato de
medição) — ele não deve ser citado como ganho.
\end{resultado}

\begin{metodo}
\textbf{Usar o perfil local:}
\begin{enumerate}[leftmargin=1.4em,itemsep=1pt]
  \item Use o arquivo \pth{opencode.json} (perfil mínimo):
         \pth{model: local/on-device}
  \item Importe os modelos afinados no Ollama antes de executar.
  \item Para a versão de nuvem/catálogo completo, use a configuração padrão
         \pth{opencode.json} (216 agentes, 7 MCPs).
\end{enumerate}
\end{metodo}

\begin{aviso}
Todos os modelos locais empataram no ``tier difícil'' (7/11 em 11 perguntas) e o
modelo \pth{Qwen3-0.6B} marcou 8/11 — diferença \textbf{não significativa} com
$n=11$. Não declarar o modelo local ``igual'' ou ``superior'' a bigger models: é
desempenho computacionalmente paralelo, sem significância estatística.
\end{aviso}

\begin{niveis}
\textbf{Intuição:} o ecossistema pode rodar ``na sua máquina'' — e um ajuste fino
deixou a conversa muito mais rápida.\\
\textbf{Implementação:} escolha de janela de contexto menor + \pth{num_ctx} limita o
bloqueio do servidor; keep-alive e limite de modelos evitam recarregas.\\
\textbf{Formalização:} nuance entre \emph{custo de contexto} e \emph{custo de inferência}:
a otimização é de decodificação sob restrição de memória, com hipótese de
equivalência não comprovada para a variedade de tarefas (amostra $n=11$).
\end{niveis}

\clearpage
\section{Executores externos e credenciais}

\elemento{CLIs externas opcionais}{\metainfo{ID}{U-04} \hfill
\metainfo{Rota}{\pth{specs/SPEC-935-R607.md}}}

\begin{descricao}
O Core pode delegar tarefas a executores externos (em vez de ``reconstruir''
capacidades). A invocação é tolerante: se a CLI não existe, avisa e segue.
\end{descricao}

\begin{resultado}
Invocáveis por skill própria (cada uma com seu grelhado de credenciais):\par
\smallskip
{\footnotesize
\begin{tabularx}{\textwidth}{lYl}
\toprule
Executor & Spec & Credencial exigida \\
\midrule
\pth{GEMINI_API_KEY} ou Vertex \\
\pth{DEEPSEEK_API_KEY}) \\
\pth{goose} & SPEC-935-R598 & conta/config da CLI Goose \\
\pth{seed-data} & SPEC-935-R599 & chave de provedor (ex.: OpenRouter) \\
\pth{haystack} & SPEC-935-R604 & depende do pipeline montado \\
NotebookLM bridge & SPEC-935-R601 & sessão \texttt{nlm\ login\ --check} + modelo Gemini \\
MiroFish proxy & SPEC-935-R605 & serviço MiroFish \\
\bottomrule
\end{tabularx}
}

\textbf{Nota:} execução com chave exige instrução explícita do operador; o Core nunca
instala esses executores por conta própria (apenas orienta).
\end{resultado}

\begin{limite}
Esses executores são \emph{fontes de opinião}, não o \emph{motor de qualidade} do
Core: o orquestrador permanece dono do ciclo SDD/TDD e qualquer saída externa só é
declarada válida após validação local (anti-overclaim R110).
\end{limite}

\begin{aviso}
Existe skill que usa \emph{APIs internas não documentadas} (bridge NotebookLM via
cookies do navegador): \textbf{pode quebrar sem aviso}. Use com clareza de risco.
\end{aviso}

\section{Verificação após mudanças}

\elemento{Ritual de conclusão de mudança}{\metainfo{ID}{U-05} \hfill
\metainfo{Camada}{5 — Governança local}}

\begin{processo}
\begin{enumerate}[leftmargin=1.4em,itemsep=1pt]
  \item Rode \texttt{python3\ -m\ marceloclaro.cli\ doctor} (ambiente saudável).
  \item Para funcionalidade nova: escreva o spec em \pth{specs/SPEC-935-R*.md} e
        os testes \pth{tests/test_r*.py}.
  \item Valide a suíte de testes da família relevante.
  \item Registre o ciclo: \pth{evolution_registry.record(...)}.
  \item Se editou cards de agentes: regenere \pth{opencode.json} com
        \texttt{python3\ -m\ integrations.opencode\_cli}.
\end{enumerate}
\end{processo}

\begin{descricao}
Este manual é a fonte didática de como operar o Core; o \pth{MANUAL.md} traz a
versão concisa executável e o \pth{ARCHITECTURE.md} o detalhe técnico.
\end{descricao}


\section{Resumo e exercícios}

\begin{resultado}
\textbf{O que você deve ter aprendido:} (1) instalar é criar um ambiente
reprodutível, não só baixar dependências; (2) o \pth{doctor} é o exame de
saúde antes de operar; (3) perfis Lean/Full/Premium são sobre memória e
soberania local.
\textbf{Exercício sugerido:} rode \texttt{python3\ -m\ marceloclaro.cli\ doctor} do
seu terminal e anote cada coluna do resultado — o que cada um verifica.
\end{resultado}

% ============================================================

%  Gerado por gen_mod14_ativacao.py (R613) — seção de ativação e cálculo
%  para o módulo mod-01c-instalacao.tex. Edições manuais são preservadas nas próximas
%  execuções (marcador impede reescrita).
% ============================================================

\section[Leitura guiada]{Leitura guiada — Instalação e ambiente: preparar, executar e conferir}

\elemento{Leitura guiada — Instalação e ambiente}{\metainfo{ID}{N-01c} \hfill \metainfo{Ciclo}{R614} \hfill \metainfo{Estilo}{calibrar antes de medir}}

\begin{descricao}
\textbf{A pergunta.} Como outra pessoa pode reconstruir o ambiente em que um
resultado foi produzido? Fixar arquivos de dependência ajuda, mas não basta:
versões do sistema, serviços externos, credenciais e fontes de aleatoriedade
também podem alterar uma execução. O comando
\texttt{python3\ -m\ marceloclaro.cli\ doctor} verifica um conjunto definido de
condições locais e informa quais passaram, falharam ou geraram avisos. É uma
primeira inspeção do ambiente; não garante que toda tarefa funcione nem que uma
saída possa ser reproduzida bit a bit.
\end{descricao}

\begin{definicao}
\textbf{Grandezas e definições.}
\begin{itemize}[leftmargin=1.3em,itemsep=1pt]
  \item \textbf{Reprodutibilidade}: registrar insumos, versões, procedimento e condições para que outra execução possa ser reconstruída e comparada. Saídas idênticas dependem também de determinismo.
  \item \textbf{Perfil de dependências}: conjuntos base, de desenvolvimento e científico (evita instalar o que não se usa).
  \item \textbf{Diagnóstico de saúde} ($h$): conjunto de verificações locais; nesta configuração, o \pth{doctor} reporta 20 checks e separa aprovações de avisos e falhas.
  \item \textbf{Custo de detecção} ($C_d$): esforço para \emph{descobrir} um defeito, que cresce quanto mais tarde ele é achado.
\end{itemize}
\end{definicao}

\begin{metodo}
\textbf{Dedução passo a passo.} Por que verificar \emph{antes} de operar?
\begin{enumerate}[leftmargin=1.4em,itemsep=2pt]
  \item Um diagnóstico inicial pode revelar dependências ausentes antes que elas interrompam uma tarefa.
  \item Se a falha aparecer durante a execução, será necessário identificar quais etapas foram afetadas e, possivelmente, repeti-las.
  \item O custo depende do tipo de falha, do estágio em que foi descoberta e do trabalho já realizado; este repositório não mede uma razão universal de custo.
  \item Portanto, rode o \pth{doctor} quando o procedimento do projeto o exigir e interprete cada aviso conforme a tarefa que pretende executar.
\end{enumerate}
\end{metodo}

\begin{resultado}
\textbf{Exemplo resolvido.} Na execução usada para esta revisão, o campo
\texttt{duration\_seconds} do diagnóstico reportou $3{,}39$ segundos. Se o mesmo
comando fosse executado antes de $12$ tarefas, a divisão simples daria
$3{,}39/12 \approx 0{,}28$ segundo de diagnóstico por tarefa. Esse cálculo só
distribui o tempo observado: não prova que todas as tarefas precisam de uma
verificação nova nem que o tempo será igual em outra máquina. \emph{Confira no seu
ambiente:} rode \texttt{python3\ -m\ marceloclaro.cli\ doctor} e use o valor de
\texttt{duration\_seconds} retornado naquela execução.
\end{resultado}

\begin{resultado}
\textbf{Sequência de conferência.} Para mudanças sujeitas ao fluxo SDD/TDD deste
repositório, consulte o diagnóstico, a especificação aplicável, os testes e o
registro de evolução; regenere \pth{opencode.json} apenas quando a alteração
envolver os catálogos que o produzem. A ordem exata depende da tarefa. Para
reproduzir o inventário de checks, leia a seção \emph{checks} da saída do
\pth{doctor}; a duração, por si só, não informa quantos checks passaram.
\end{resultado}

\begin{aviso}
\textbf{Limite de validade.} O \pth{doctor} testa \emph{o ambiente}, não o
\emph{mérito} das tarefas. Ambiente verde significa ``pronto para operar'', nunca
``resultado verificado''.
\end{aviso}

\begin{resultado}
\textbf{Exercícios.} (1) Descreva o que muda entre os perfis base, dev e
científico. (2) Por que toda referência deste livro exige ``data de acesso''
(Módulo 11)?
\end{resultado}

% ATIVACAO-R613
\section[Ativação e cálculo]{Ativação e cálculo — Instalação e ambiente}
\elemento{ativ-01c}{\metainfo{Ciclo}{R613} \hfill \metainfo{Módulo}{mod-01c-instalacao.tex}}

\begin{processo}
GATILHO. ``instalar'', ``ambiente mínimo'', ``doctor''.
ROTA. linha de comando; gate: \texttt{python3\ -m\ marceloclaro.cli\ doctor} com checks\_failed = 0.
FUNÇÃO. garantir ambiente reprodutível antes de qualquer ativação de agente ou skill.
\end{processo}

\begin{resultado}
CÁLCULO. saúde do ambiente = 1 se checks\_failed = 0, senão 0; o manual não segue com 0.
\end{resultado}

\begin{descricao}
JUSTIFICATIVA TEÓRICA. alegações computacionais exigem padrão mínimo de reprodução do ambiente. O lastro teórico vem da citação que segue: recorte conferido em 29/09/2026, com a página de origem e tradução livre e fiel.
\end{descricao}

\begin{aviso}
LIMITE E ANTI-OVERCLAIM. ambiente verde não implica qualidade de resultado — apenas reprodutibilidade do processo. A verificação atesta a existência e a
localização da fonte (R110), não o mérito da afirmação.
\end{aviso}

\noindent\textit{Interpretação:} A instalação reproduzível depende de condições conjuntas: instruções executáveis, versões identificadas, entradas conhecidas e resultado observável. Como verificação didática, escrevemos $R=I\cdot V\cdot D\cdot S$, em que cada indicador vale $1$ quando a condição foi registrada e conferida, e $0$ caso contrário. Se uma versão não é fixada, por exemplo, $V=0$ e o produto é zero; isso não prova que o programa falhará, mas mostra que a execução ainda não pode ser repetida sob as mesmas premissas \citref{PENG, R. D. Reproducible research in computational science. \emph{Science}, v.~334, n.~6060, p.~1226-1227, 2011}{10.1126/science.1213847}{p.~1226--1227 (resumo, p.~1226). ancora o gate de instalação: reprodutibilidade é o padrão mínimo para julgar alegações quando a replicação integral não é possível — o que o doctor verifica antes de todo pipeline.}{Reproducibility has the potential to serve as a minimum standard for judging scientific claims when full independent replication of a study is not possible}{A reprodutibilidade tem o potencial de servir como padrão mínimo para julgar alegações científicas quando a replicação independente completa de um estudo não é possível.}{PENG, 2011, p.~1226}
